% =============================================================================
% Template Pandoc LaTeX pour la classe memoirerecherche
% Utilisé par bookdown::pdf_book via rapport.Rmd
% =============================================================================

% Pré-déclarer l'option [table] pour xcolor AVANT le chargement de la classe.
% memoirerecherche.cls fait \RequirePackage[table]{xcolor} ; si kableExtra
% (ou un autre paquet) tente ensuite \usepackage{xcolor}, LaTeX peut
% signaler un « Option clash for package xcolor ». La ligne suivante
% garantit la compatibilité dans tous les cas.
\PassOptionsToPackage{table}{xcolor}

\documentclass[11pt,a4paper]{memoirerecherche}

% --- Géométrie ----------------------------------------------------------------

% --- Encodage et polices (XeLaTeX) --------------------------------------------
\usepackage{fontspec}
\defaultfontfeatures{Ligatures=TeX,Scale=MatchLowercase}

% --- Langue -------------------------------------------------------------------
\usepackage[french]{babel}
\usepackage{csquotes}

% --- Mathématiques ------------------------------------------------------------
\usepackage{amsmath,amssymb}

% --- Tableaux et flottants ----------------------------------------------------
% Note : graphicx, xcolor, colortbl, hyperref, caption déjà chargés par memoirerecherche.cls
\usepackage{longtable}
\usepackage{booktabs}
\usepackage{array}
\usepackage{calc}   % Requis par Pandoc ≥ 3.x pour \real{} dans les largeurs de colonnes
\usepackage{multirow}
\usepackage{float}

% --- Graphiques (maxwidth/maxheight pour Pandoc) -------------------------------
\makeatletter
\def\maxwidth{\ifdim\Gin@nat@width>\linewidth\linewidth\else\Gin@nat@width\fi}
\def\maxheight{\ifdim\Gin@nat@height>\textheight\textheight\else\Gin@nat@height\fi}
\makeatother
\setkeys{Gin}{width=\maxwidth,height=\maxheight,keepaspectratio}

% --- Listes -------------------------------------------------------------------
\providecommand{\tightlist}{%
  \setlength{\itemsep}{0pt}\setlength{\parskip}{0pt}}

% --- Configuration hyperref (package déjà chargé par la classe) ---------------
\hypersetup{
  pdftitle={L'algorithme Affinity Propagation},
  colorlinks=true,
  linkcolor=urcabrown,
  citecolor=urcabrown,
  urlcolor=urcalightbrown,
  breaklinks=true
}
\urlstyle{same}

% --- En-têtes / pieds de page -------------------------------------------------
\usepackage{fancyhdr}
\pagestyle{fancy}
\fancyhf{}
\fancyhead[L]{\small\leftmark}
\fancyhead[R]{\small\thepage}
\renewcommand{\headrulewidth}{0.4pt}

% --- Bibliographie (biblatex) -------------------------------------------------
\usepackage[]{biblatex}
\addbibresource{bibliography.bib}

% --- Comptage automatique des références bibliographiques ---------------------
% Le compteur citnum (défini dans memoirerecherche.cls) est incrémenté pour
% chaque entrée imprimée dans la bibliographie ; \total{citnum} sur la page
% de titre affichera la valeur finale (disponible dès le second passage LaTeX).
\newcounter{citnum}
\AtEveryBibitem{\stepcounter{citnum}}

% --- Coloration syntaxique Pandoc ---------------------------------------------
% Pré-définir l'environnement Shaded (requis par les macros de coloration
% Pandoc qui utilisent \renewenvironment) et la couleur shadecolor.
\usepackage{tcolorbox}
\usepackage{framed}
\definecolor{shadecolor}{RGB}{248,248,248}
\newenvironment{Shaded}{\begin{snugshade}}{\end{snugshade}}

% --- Header-includes depuis le YAML -------------------------------------------
\makeatletter
\@ifpackageloaded{caption}{}{\usepackage{caption}}
\AtBeginDocument{%
\ifdefined\contentsname
  \renewcommand*\contentsname{Table of contents}
\else
  \newcommand\contentsname{Table of contents}
\fi
\ifdefined\listfigurename
  \renewcommand*\listfigurename{List of Figures}
\else
  \newcommand\listfigurename{List of Figures}
\fi
\ifdefined\listtablename
  \renewcommand*\listtablename{List of Tables}
\else
  \newcommand\listtablename{List of Tables}
\fi
\ifdefined\figurename
  \renewcommand*\figurename{Figure}
\else
  \newcommand\figurename{Figure}
\fi
\ifdefined\tablename
  \renewcommand*\tablename{Table}
\else
  \newcommand\tablename{Table}
\fi
}
\@ifpackageloaded{float}{}{\usepackage{float}}
\floatstyle{ruled}
\@ifundefined{c@chapter}{\newfloat{codelisting}{h}{lop}}{\newfloat{codelisting}{h}{lop}[chapter]}
\floatname{codelisting}{Listing}
\newcommand*\listoflistings{\listof{codelisting}{List of Listings}}
\makeatother
\makeatletter
\makeatother
\makeatletter
\@ifpackageloaded{caption}{}{\usepackage{caption}}
\@ifpackageloaded{subcaption}{}{\usepackage{subcaption}}
\makeatother

% --- Variables personnalisées → commandes memoirerecherche.cls ----------------
\diplome{À compléter}
\module{Apprentissage non supervisé}
\enseig{Frédéric Blanchard, Philippe Regnault}
\evaluation{Rapport}
\coefficient{À compléter}
\anac{2026-2027}

% Résumé
\resume{}

% Mots-clés
\keywords{}

% Dépôt Git et base de données
\gitrepo{https://gitlab-mi.univ-reims.fr/dea0014/affinity\_propagation}
\database{}

% Logo de la page de titre
\logo{logoSESGURCA.jpg}

% Métriques de la page de titre — comptage automatique
% Les compteurs figure, table et citnum sont gérés automatiquement par
% totcount (voir memoirerecherche.cls). Les lignes ci-dessous ne sont
% conservées que pour ne pas casser un éventuel code qui référencerait
% \@totfig, \@tottab, etc.
\totfig{--}\tottab{--}\totref{--}\tottheo{0}

% --- Titre, auteur, date ------------------------------------------------------
\title{L'algorithme Affinity Propagation}
% N.B. : ne PAS utiliser \and ici — la classe memoirerecherche.cls affiche
% \@author dans un minipage/tabular ; or, \and est défini par article.cls
% comme \end{tabular}\begin{tabular}, ce qui casse l'imbrication de tables
% du \maketitle de la classe et provoque « ! Missing } inserted. ».
\author{Hervé ADJAKOSSA \& Marie DE AMORIM PEREIRA DA SILVA}
\date{2026-11-05}

% =============================================================================
\begin{document}

\maketitle
\newpage

% --- Table des matières -------------------------------------------------------
{
\setcounter{tocdepth}{3}
\tableofcontents
\newpage
}

% --- Liste des tables ---------------------------------------------------------

% --- Liste des figures --------------------------------------------------------

% --- Corps du document --------------------------------------------------------
\newpage

\section{INTRODUCTION}\label{introduction}

Dans le cadre de la matière apprentissage non supervisé, enseignée par
Monsieur Frédéric BLANCHARD et Monsieur Philippe REGNAULT, nous étudions
l'algorithme ``Affinity Propagation''.

Pour commencer, il semble important d'établir une première distinction
entre l'apprentissage supervisé et non supervisé. L'apprentissage
non-supervisé, aussi appelé apprentissage automatique, consiste à
utiliser des algorithmes pour apprendre des structures, des relations ou
des règles à partir de données. Au lieu de programmer manuellement
toutes les règles que l'ordinateur doit suivre, on lui fournit des
données et une méthode qui lui permet d'en extraire certaines
informations. L'apprentissage supervisé en revenche, consiste à utiliser
des algorithmes pour apprendre des relations entre des données et des
résultats connus, afin de prédire ces résultats pour de nouvelles
observations. La différence majeure réside donc dans la connaissance (ou
non) du résultat à prédire sur les données d'entraînement. Ce projet
s'inscrit donc dans une démarche d'apprentissage non-supervisé.

Le terme d'algorithme a été mentionné à plusieurs reprises, il nous
semble donc important de commencer par le définir. Un algorithme est un
processus systématique de résolution d'un problème. Un algorithme
d'apprentissage automatique est plus élaboré : il utilise les données
pour construire une solution à un problème donné. Il existe de
nombreuses méthodes, car tous les problèmes ne se ressemblent pas.
Certaines méthodes cherchent à prédire une valeur, d'autres à classer
des observations, d'autres encore à découvrir des groupes ou des
relations cachées dans les données.

Nous nous intéresserons donc, dans ce rapport, à l'algorithme ``Affinity
Propagation''. Notre travail permettra donc notamment de comprendre le
but de cet algorithme, le comparer à d'autres algorithmes ayant un but
similaire afin d'identifier les situations les plus propices à son
utilisation, puis de comprendre les fondements mathématiques
sous-jacents. Finalement, nous préciserons une implémentation existante
sour R avant de mettre en place une application permettant de se
familiariser avec Affinity Propagation.

\newpage

\section{I. POSITIONNEMENT DE L'ALGORITHME AFFINITY PROPAGATION DANS
L'APPRENTISSAGE
AUTOMATIQUE}\label{i.-positionnement-de-lalgorithme-affinity-propagation-dans-lapprentissage-automatique}

\subsection{A. La classification non supervisée : le
clustering}\label{a.-la-classification-non-supervisuxe9e-le-clustering}

L'apprentissage non supervisé regroupe plusieurs familles de méthodes
ayant pour objectif d'extraire des informations à partir de données non
étiquetées. Parmi celles-ci, le \emph{clustering}, également appelé
classification non supervisée, consiste à répartir les observations d'un
jeu de données en plusieurs groupes, appelés clusters.

L'objectif est de réunir au sein d'un même groupe les observations
présentant des caractéristiques similaires, tout en distinguant autant
que possible les observations appartenant à des groupes différents.
Contrairement à la classification supervisée, les groupes ne sont pas
définis à l'avance. L'algorithme doit donc identifier lui-même une
organisation pertinente des données, à partir des informations dont il
dispose.

Le clustering peut répondre à différents objectifs, notamment identifier
des profils similaires au sein d'une population, regrouper des individus
présentant des caractéristiques communes ou encore découvrir des
structures dans des données complexes.

Il existe cependant plusieurs manières de construire ces groupes. Les
méthodes de clustering se distinguent notamment par la façon dont elles
mesurent les ressemblances entre observations et par les règles
utilisées pour former les groupes.

\subsection{B. Les principales familles de méthodes de
clustering}\label{b.-les-principales-familles-de-muxe9thodes-de-clustering}

Il existe différentes approches pour réaliser une classification non
supervisée. Parmi les plus connues, on peut notamment distinguer les
méthodes de partitionnement, les méthodes hiérarchiques et les méthodes
fondées sur la densité.

Concernant les méthodes de partitionnement,elles cherchent à répartir
directement les observations en un nombre déterminé de groupes.
L'algorithme K-means en est un exemple classique. Il consiste à
attribuer chaque observation au centre de groupe le plus proche, puis à
actualiser ces centres jusqu'à obtenir une certaine stabilité. Cette
approche nécessite généralement de définir le nombre de groupes souhaité
avant de lancer l'algorithme. On verra donc qu'Affinity Propagation est
un algorithme particulier de cette famille.

Les méthodes hiérarchiques quant à elles construisent une organisation
des observations sous la forme d'une hiérarchie de groupes. Dans une
approche ascendante, chaque observation constitue initialement un groupe
individuel, puis les groupes les plus proches sont progressivement
fusionnés. Les résultats peuvent être représentés à l'aide d'un
dendrogramme, qui permet de visualiser les différentes étapes de
regroupement.

Finalement, les méthodes fondées sur la densité cherchent à identifier
des zones dans lesquelles les observations sont particulièrement
concentrées. DBSCAN (Density-Based Spatial Clustering of Applications
with Noise) est un exemple de cette famille. Il permet notamment
d'identifier des groupes de formes variées et de distinguer les
observations considérées comme du bruit.

Ces différentes approches reposent sur des principes distincts et
présentent chacune des avantages et des limites selon la structure des
données étudiées.

\subsection{C. La place de l'Affinity
Propagation}\label{c.-la-place-de-laffinity-propagation}

L'algorithme Affinity Propagation, proposé par Brendan J. Frey et
Delbert Dueck en 2007 \autocite{frey2007clustering}, constitue une
approche particulière du clustering.

Contrairement à K-means, il ne nécessite pas de fixer directement le
nombre de groupes avant l'exécution de l'algorithme. Il repose sur les
similarités entre les observations et cherche à identifier des
observations représentatives, appelées ``exemplars'', qui serviront de
représentants aux différents groupes. Chaque observation est ensuite
associée à l'un de ces exemplaires, ce qui permet de constituer les
clusters.

Le fonctionnement de l'algorithme repose sur des échanges itératifs de
messages entre les observations. Ces messages permettent progressivement
de déterminer quelles observations sont les plus pertinentes pour
représenter les groupes.

L'Affinity Propagation se distingue ainsi par son principe de sélection
des représentants : plutôt que de construire des centres artificiels,
comme dans K-means, il sélectionne des observations réellement présentes
dans le jeu de données. Cette particularité constitue l'un des éléments
centraux de son fonctionnement et sera étudiée plus en détail dans les
parties suivantes.

% --- Bibliographie ------------------------------------------------------------
\printbibliography[heading=bibintoc,title={Bibliographie}]

\end{document}
